\documentclass[11pt]{report}

% ============================================
% IA Generativa para Data Science — Notas de Estudo
% Consolidado a partir dos Jupyter notebooks
% ============================================

% --- Pacotes ---
\usepackage[utf8]{inputenc}
\usepackage[T1]{fontenc}
\usepackage[brazil]{babel}
\usepackage{amsmath,amssymb}
\usepackage{graphicx}
\usepackage{hyperref}
\usepackage{booktabs}
\usepackage{listings}
\usepackage{xcolor}

% --- Configurações ---
\graphicspath{{figures/}}

\lstset{
    basicstyle=\ttfamily\small,
    breaklines=true,
    frame=single,
    numbers=left,
    language=Python
}

% ============================================
\begin{document}

\title{Inteligência Artificial Generativa\\ {\it From Deep Learning to Large Language Models}}
\author{Lucas A. Souza}
\date{2026}



\maketitle
\begin{abstract}
    Este documento apresenta um resumo das principais ideias e conceitos da Inteligência Artificial Generativa, desde os fundamentos do Deep Learning até os modelos de linguagem (LLMs). A abordagem é baseada em estudos e implementações práticas.
\end{abstract}


\tableofcontents

% ============================================
% PARTE I — FUNDAMENTOS
% ============================================

\part{Fundamentos}

\chapter{Fundamentos de Deep Learning}
% Fundamentos de Deep Learning

Perceptron, MLP, forward pass, backpropagation, funções de ativação, funções de perda, gradiente descendente, Adam, regularização, Dropout, BatchNorm e validação.

% Notas detalhadas e código nos notebooks em contents/01/deep_learning/


\chapter{NLP Clássico e Representação de Texto}
% NLP Clássico e Representação de Texto

Tokenização, n-grams, vetores de palavras (word2vec, GloVe), classificação de texto, embeddings, TF-IDF, bag of words e métricas de similaridade.

% Notas detalhadas e código nos notebooks em contents/02/nlp/


\chapter{RNN, LSTM e GRU}
% RNN, LSTM e GRU

Redes recorrentes, vanishing gradient, LSTM, GRU, modelagem de sequências, geração e classificação de texto. Limitações frente aos Transformers.

% Notas detalhadas e código nos notebooks em contents/03/rnn_lstm/


% ============================================
% PARTE II — TRANSFORMERS E LLMs
% ============================================

\part{Transformers e Large Language Models}

\chapter{Attention e Transformers}
% Attention e Transformers

Query, Key, Value, self-attention, masked attention, cross-attention, multi-head attention, positional encoding, encoder, decoder.

% Notas detalhadas e código nos notebooks em contents/04/transformers/


\chapter{Large Language Models}
% Large Language Models

Transformers, tokenização, modelos pré-treinados, BERT, GPT, fine-tuning, prompting, RLHF e alinhamento.

% Notas detalhadas e código nos notebooks em contents/05/llm/


% ============================================
% PARTE III — RAG E FRAMEWORKS
% ============================================

\part{RAG, Frameworks e Agentes}

\chapter{Embeddings, Busca Vetorial e RAG}
% Embeddings, Busca Vetorial e RAG

Loaders, chunking, embeddings, vector stores, retrieval, reranking, avaliação de RAG, chat com documentos.

% Notas detalhadas e código nos notebooks em contents/06/rag/


\chapter{LangChain, LangGraph e LlamaIndex}
% LangChain, LangGraph e LlamaIndex

Models, prompts, parsers, chains, memória, agentes, tools, workflows, ingestão, indexação, query engines.

% Notas detalhadas e código nos notebooks em contents/07/langain/


\chapter{Agentes de IA}
% Agentes de IA

Reasoning loop, tools, actions, function calling, estados, transições, memória, human-in-the-loop, avaliação de agentes.

% Notas detalhadas e código nos notebooks em contents/08/agents/


\chapter{Model Context Protocol (MCP)}
% Model Context Protocol (MCP)

Arquitetura MCP, clientes, servidores, tools, resources, prompts, SDKs e integração com agentes.

% Notas detalhadas e código nos notebooks em contents/09/mcp/


% ============================================
% PARTE IV — PRODUÇÃO
% ============================================

\part{Produção e Deploy}

\chapter{OpenAI API e Azure OpenAI}
% OpenAI API e Azure OpenAI

SDK, Responses API, streaming, tools, structured outputs, RAG com Azure AI Search, busca híbrida e ranking semântico.

% Notas detalhadas e código nos notebooks em contents/10/openai_api/


\chapter{FastAPI}
% FastAPI

Rotas, parâmetros, Pydantic, validação, dependências, segurança, autenticação, testes e aplicações modulares.

% Notas detalhadas e código nos notebooks em contents/11/fastapi/


\chapter{Docker e Deploy}
% Docker e Deploy

Containers, imagens, Dockerfile, volumes, redes, Docker Compose, orquestração de API e banco vetorial.

% Notas detalhadas e código nos notebooks em contents/12/docker/


\chapter{Observabilidade, Avaliação e Guardrails}
% Observabilidade, Avaliação e Guardrails

Tracing, latência, erros, tokens, custo, guardrails, suítes de avaliação, cache, fallback, LLMOps.

% Notas detalhadas e código nos notebooks em contents/13/observabilidade/


% ============================================
\end{document}
